\documentclass[12pt,a4paper]{article}

\usepackage[brazil]{babel}
\usepackage[utf8]{inputenc}
\usepackage[T1]{fontenc}
\usepackage{lmodern}
\usepackage{amsmath, amssymb, amsthm}
\usepackage{graphicx}
% Os PDFs de figs/ são regeneráveis e não-versionados; antes de compilar:
%   python scripts/gerar_figs_paper.py
\graphicspath{{figs/}}
\usepackage{booktabs}
\usepackage{hyperref}
\usepackage[round]{natbib}
\usepackage{geometry}
\geometry{margin=1in}

\newtheorem{proposition}{Proposição}

\title{Leniency Conditional on Critical Mass:\\
A Conditional-Deposit Channel for Unilateral Conduct in Digital Markets}
\author{L.\thanks{O autor mantém o repositório \texttt{waas-antitrust} de forma independente; as opiniões aqui expressas não vinculam qualquer instituição.}}
\date{\today}

\begin{document}
\maketitle

\begin{center}\textbf{Rascunho de trabalho (v0.2.0) --- não citar.}\end{center}

\begin{abstract}
Os programas clássicos de leniência, ao explorarem a instabilidade interna do dilema do prisioneiro entre conspiradores externos, perdem tração diante das condutas predominantemente unilaterais que caracterizam o enforcement antitruste de mercados digitais \citep{cremer2019competition}. Propomos a \emph{Leniência Condicionada à Massa Crítica} (LCMC), um \emph{canal de depósito condicional} operado pelo CADE que implementa o desenho de \emph{information escrow} de \citet{ayres2012information} no domínio antitruste, com análogo operacional em \citet{callisto}: depósitos individuais de denúncia permanecem selados até que um gatilho de massa crítica $q_{\min}\cdot n$ sobre a mesma firma seja atingido, hipótese em que se abrem simultaneamente para o procedimento administrativo. A base normativa do canal é autônoma --- Art.~4º, II e III, da Lei nº~12.529/2011 \citep{lei12529} combinado com a Lei nº~9.784/1999 \citep{lei9784} ---, dispensando nova lei e dissociando a coordenação do incentivo monetário. Ao canal podem ser acoplados, de modo facultativo e ortogonal, cinco instrumentos de internalização: canal puro sem recompensa; recompensa via Termo de Compromisso de Cessação (TCC) sob o Art.~12 da Resolução CADE nº~21/2018 \citep{res21cade}; saída coordenada com vesting acelerado \emph{à la} Hirschman; crédito tributário; e leniência criminal individual. Quando a recompensa via TCC --- o instrumento \emph{Whistleblower-as-a-Service} (WaaS) --- é o acoplamento ativo, a firma maximiza a margem $D - W$, e a fronteira IC-F* ($D > W$) é satisfazível na jurisprudência atual do Art.~12 (Regime~B, infralegal; construção controvertível, sujeita a anulação judicial). A análise empírica usa modelo baseado em agentes em Mesa~3.x com três populações (trabalhador, empresa, autoridade) e seis fases por tique. A calibração formal por Nelder-Mead em (fração de violadoras, taxa de capacidade) $=(0{,}323;\ 0{,}481)$ entrega erro relativo de 6,65\% sobre o alvo operacional de TCCs/ano, dentro do intervalo de confiança \emph{bootstrap} de 95\%, e implica universo CADE de $N^\star \approx 1\,679$ firmas como predição falsificável. A validação institucional externa apoia-se em dois marcos contemporâneos: o programa de recompensa antitruste do DOJ-ATR \citep{dojatr2025}, instituído em julho de 2025 por ato infralegal em parceria com o USPS, oferece paralelo direto à via administrativa do Regime~B; a Whistleblower Tool da DMA \citep{dmatool2024}, lançada em abril de 2024 sem recompensa associada, oferece contrafactual para o canal puro. Reconhecem-se duas limitações: a forma forte da Proposição~5 (erosão Coleman do substrato cooperativo) foi refutada por varredura multi-seed dedicada, sobrevivendo apenas a forma fraca; e a unicidade do equilíbrio de coordenação (Proposição~2) sob heterogeneidade de arquétipos e papéis segue conjectura aberta.
\end{abstract}

\section{Introdução}

Programas de leniência reconfiguraram o enforcement antitruste das últimas três décadas ao oferecer imunidade qualificada ao primeiro membro de um cartel que coopere com a autoridade. A análise formal de \citet{spagnolo2004divide} mostrou que descontos progressivos por posição na fila sustentam, em equilíbrio bayesiano-perfeito, a delação como resposta dominante em cartéis de população finita; \citet{harrington2015when} explicitaram, em sequência, as condições agregadas sob as quais um programa de leniência reduz a prevalência cartelizada no longo prazo. A evidência empírica em jurisdições maduras --- DOJ Antitrust Division desde 1993, Comissão Europeia desde 1996, CADE desde 2000 --- corroborou que o instrumento opera por desestabilização interna de uma coordenação que já existe entre concorrentes externos.

O sucesso depende, porém, da existência de um \emph{cartel}: conduta coordenada entre firmas distintas cuja instabilidade interna a leniência explora. \citet{cremer2019competition} documentam que parte central das condutas potencialmente anticompetitivas em mercados digitais é estruturalmente distinta --- autopreferência de busca, restrições \emph{anti-steering} sobre desenvolvedores, recusa de acesso a dados, aquisições assassinas (\emph{killer acquisitions}), exclusividade em \emph{marketplace}. São condutas \emph{unilaterais}, sem co-conspirador externo a quem se possa oferecer imunidade; a informação relevante reside, tipicamente, no interior de pequenos núcleos técnicos da própria firma (engenharia de \emph{ranking}, \emph{ad-tech}, \emph{corp dev}). A leniência clássica, desenhada para o dilema do prisioneiro entre firmas, perde tração nesse domínio: não há a quem delatar, há apenas a quem testemunhar.

O problema central, sustentamos, não é a punição da firma --- já endereçada pelo arsenal sancionatório da Lei nº~12.529/2011 --- nem o pagamento ao denunciante individual --- já contemplado, sob arquitetura distinta, pela Lei nº~13.608/2018 \citep{lei13608}. O problema é a \emph{coordenação intra-firma}: o trabalhador que detém informação relevante não se identifica isoladamente, porque o custo individual de represália suportado em isolamento é proibitivo; o equilíbrio de silêncio que daí resulta é o caso paradigmático da \emph{lógica da ação coletiva} de \citet{olson1965logic}, do modelo de \emph{thresholds} de \citet{granovetter1978threshold} e do \emph{contágio complexo} de \citet{centola2007complex} sobre redes intra-organizacionais.

Este artigo propõe a \emph{Leniência Condicionada à Massa Crítica} (LCMC) como mecanismo de coordenação adequado ao domínio. O núcleo do desenho é um \emph{canal de depósito condicional} operado pelo CADE: trabalhadores depositam denúncias seladas em escrow institucional, à maneira do \emph{information escrow} de \citet{ayres2012information}; os depósitos só se abrem simultaneamente quando um gatilho de massa crítica $q_{\min}\cdot n$ sobre a mesma firma é atingido, replicando \emph{mutatis mutandis} o desenho \emph{all-or-nothing} de \citet{callisto} para denúncias condicionais de assédio em campus universitário. O CADE figura como custodiante neutro do escrow; nenhum depósito individual permanece exposto isoladamente. A base normativa do canal é autônoma no Art.~4º, II e III, da Lei nº~12.529/2011 \citep{lei12529} combinado com a Lei nº~9.784/1999 \citep{lei9784} sobre processo administrativo federal: o canal é implementável por ato infralegal do CADE, sem nova lei.

Ao canal podem ser acoplados, de modo opcional e ortogonal, cinco instrumentos de internalização: o canal puro (sem instrumento acoplado, equivalente operacional da DMA Tool); a recompensa via TCC sob re-caracterização do Art.~12 da Resolução CADE nº~21/2018 \citep{res21cade} --- o instrumento \emph{Whistleblower-as-a-Service} (WaaS) que dá nome à literatura precedente do projeto e que satisfaz a aritmética IC-F* ($D > W$) sobre a firma, sujeita à validação judicial; o vesting acelerado de Hirschman como saída coordenada da firma; o crédito tributário (sub-regime $C_\varphi$); e a leniência criminal individual (sub-regime $C_p$). O acoplamento é, em todo caso, lógica e juridicamente separável do canal. A Seção~\ref{sec:mecanismo} detalha o mecanismo e a charneira normativa de cada regime; a Seção~\ref{sec:limitacoes}, o modelo baseado em agentes e a calibração formal; as seções subsequentes desenvolvem o enquadramento cibernético, a pré-condição sociológica e a tradução política do desenho ao contexto institucional brasileiro 2024--2027.

\section{Falha estrutural da leniência clássica em mercados digitais}

A leniência clássica é uma tecnologia de ruptura de cartéis: explora a instabilidade
interna do dilema do prisioneiro entre conspiradores externos. Mas as condutas
identificadas como prioritárias em mercados digitais \citep{cremer2019competition}
--- autopreferência, restrições verticais sobre desenvolvedores, recusa de acesso a
dados, predação algorítmica, vinculação --- são predominantemente unilaterais. Não
há co-conspirador externo a quem oferecer imunidade. A informação relevante reside
em pequenos núcleos técnicos da firma (engenharia de produto, ranking, ad-tech,
acquisitions), tipicamente 2--3 papéis primários. O número absoluto de pessoas com
conhecimento qualificado é pequeno; a coordenação \emph{intra-firma} é o problema,
não a coordenação \emph{inter-firma}.

A intuição-chave do canal de depósito condicional é trocar de jogo: se a
coordenação simultânea não emerge espontaneamente, o canal a fabrica como
serviço público --- assim como Kickstarter fabrica financiamento condicional
\emph{all-or-nothing} para projetos que jamais arrancariam de doadores isolados.
Marcos contemporâneos confirmam a transponibilidade do desenho: o
\citet{dmatool2024} oferece canal anônimo, sem recompensa, para violações do
Regulamento (UE) 2022/1925; a \citet{directive20191937} fornece proteção horizontal
anti-represália; o \citet{dojatr2025} institui o primeiro programa federal
antitruste de recompensa, 15--30\% sobre multas $\geq$~US\$~1~milhão, com primeiro
prêmio pago em janeiro de 2026 --- precedente direto para o caminho
infralegal-recompensado análogo ao Regime~B.

\section{Blocos construtivos de desenho de mecanismos}

O desenho da LCMC combina três blocos de literatura distintos. Primeiro, a teoria
de leniência ótima de \citet{spagnolo2004divide}, que mostra como descontos
progressivos por posição na fila sustentam delação em equilíbrio bayesiano-perfeito
em cartéis com população finita; transpomos o gradiente de descontos para a fila
intra-firma calibrada empiricamente por \citet{saito2021tcc} (1ª posição: 43,43\%
de desconto médio na Superintendência Geral do CADE; 2ª: 34,51\%; 3ª: 20,22\%; piso
de 15\% a partir da 9ª). Segundo, os blocos analíticos de
\citet{aubert2006impact}, que formalizam a coexistência de programas de leniência
com recompensa direta ao denunciante e identificam condições suficientes para
complementaridade entre os dois instrumentos; emprestamos a estrutura do payoff do
denunciante interno mas substituímos a noção de delação unilateral pela de
\emph{depósito condicional}. Terceiro, o template institucional do Dodd-Frank \S922
\citep{doddfrank922} --- que estabeleceu pagamento de 10--30\% sobre multas
SEC superiores a US\$~1~milhão e gerou jurisprudência judicial robusta sobre
re-caracterização do pagamento como ressarcimento à entidade lesada
(\emph{relator-style} sob a False Claims Act é o paralelo doutrinário). A
transposição ao antitruste digital exige adaptações específicas: a SEC opera sob
\emph{taxing power} federal sem reserva penal, enquanto o CADE atravessa três
competências constitucionais brasileiras (concorrencial, tributária quando ativada
$C_\varphi$, e penal quando ativada $C_p$), cada uma com sua reserva --- mapa institucional
detalhado em \S5.

A novidade não está nos blocos individualmente, mas no seu acoplamento sob
um \emph{escrow institucional}: a fila ordenada \emph{à la} Spagnolo só dispara
quando a massa crítica \emph{à la} \citet{ayres2012information} é atingida no
escrow. As recompensas \emph{à la} Aubert-Rey-Kovacic decaem com a posição via
gradiente Saito, mas o pagamento total à $k$-ésima posição só existe se
$k \geq q_{\min}\cdot n$. Esta combinação --- fila condicional com gatilho de
massa crítica --- é o que distingue a LCMC dos programas existentes.

\section{Análise comparada de programas existentes}

A Tabela~\ref{tab:comparada} sintetiza cinco programas contemporâneos ao longo
de três eixos: presença de \emph{canal} qualificado de recepção; presença de
\emph{recompensa monetária}; base normativa do instrumento (lei estatutária,
regulação infralegal, ou diretiva supranacional).

\begin{table}[t]
\centering
\small
\caption{Cinco programas contemporâneos de delação institucional.}
\label{tab:comparada}
\begin{tabular}{lccl}
\toprule
Programa & Canal & Recompensa & Base normativa \\
\midrule
SEC Whistleblower (EUA, 2010) & \checkmark & 10--30\% & Dodd-Frank \S922 \citep{doddfrank922} \\
DOJ-ATR Rewards (EUA, 2025) & \checkmark & 15--30\% & Resolução administrativa \citep{dojatr2025} \\
DMA Whistleblower Tool (UE, 2024) & \checkmark & --- & Reg.~(UE) 2022/1925 \citep{dmatool2024} \\
Diretiva (UE) 2019/1937 & --- & --- & Diretiva supranacional \citep{directive20191937} \\
CADE/leniência (BR) & --- & --- (imunidade) & Lei 12.529 Art.~86 \\
\midrule
\textbf{LCMC proposta (BR)} & \textbf{\checkmark\ (escrow)} & \textbf{opcional} & \textbf{Resolução CADE (Art.~4º Lei 12.529 + Lei 9.784)} \\
\bottomrule
\end{tabular}
\end{table}

Quatro observações estratégicas decorrem da tabela. (i) Apenas dois programas
oferecem ambos canal qualificado \emph{e} recompensa monetária --- SEC e DOJ-ATR ---
e ambos operam em jurisdição com \emph{taxing power} federal robusto;
\citet{lei12529} brasileira é estruturalmente análoga à DMA em escopo, mas
distinta em mecanismo (BR não tem canal anônimo equivalente). (ii) A
arquitetura UE separa explicitamente proteção (Diretiva 2019/1937, horizontal) e
canal (DMA Tool, setorial) --- arquitetura que preserva proteção anti-represália
mesmo onde recompensa não é institucionalmente viável.
(iii) O programa DOJ-ATR --- instituído em julho de 2025 sob a forma de
Resolução administrativa em parceria com o USPS, sem necessidade de novo
texto estatutário --- é o paralelo institucional mais próximo do Regime~B
brasileiro proposto: ambos operam o canal por ato infralegal sob autoridade
preexistente. (iv) Nenhum programa atual combina canal de depósito condicional
com massa crítica \emph{intra-firma} como gatilho; o desenho do Callisto
universitário \citep{callisto} é o análogo operacional mais próximo, mas seu
domínio (assédio sexual em campus) é estruturalmente distinto do antitruste.

A LCMC ocupa, portanto, um espaço ainda vazio na taxonomia: canal de depósito
condicional com massa crítica intra-firma, opcionalmente acoplado a recompensa
via TCC, sob via infralegal. Esse posicionamento permite implementação no
horizonte regulatório brasileiro 2024--2027 sem dependência de nova lei
federal (analisada em detalhe em \S\ref{sec:traducao_politica}).

\section{O mecanismo LCMC}
\label{sec:mecanismo}

O WaaS opera em cinco fases sequenciais. (P1)~\emph{Sinalização}: cada empregado que
observa indício de violação recebe um sinal privado ruidoso e decide, segundo seu
arquétipo \citep{hokamp2010income}, se reporta. (P2)~\emph{Depósito condicional}: o
empregado que decide reportar não envia uma denúncia tradicional --- deposita uma
denúncia condicional no escrow operado pela autoridade, \emph{à la}
\citet{ayres2012information}; o depósito é selado, e a firma não é notificada.
(P2.5)~\emph{Abertura simultânea}: o canal verifica, a cada ciclo, se a fração de
co-depósitos sobre a mesma firma atingiu o gatilho de massa crítica
$q_{\min}\cdot n$; se sim, todos os depósitos da firma se abrem simultaneamente
(\emph{all-or-nothing}), colapsando em um único caso de prova qualificada.
(P3)~\emph{Decisão da firma}: aberto o caso, sob o acoplamento monetário, a firma
compara o desconto $D$ com a recompensa total $W$ e paga quando IC-F* ($D > W$) é
satisfeita. (P4)~\emph{Autoridade}: o caso é instaurado, sujeito a uma restrição de
capacidade $\kappa$ no espírito de \citet{harrington2015when} e a uma acurácia que
cresce com a qualidade da prova. A dinâmica de coordenação subjacente ao gatilho de
P2.5 é \emph{inspirada} no jogo global de \citet{morris1998unique} e aproxima o
contágio complexo de \citet{centola2007complex} sobre redes de mundo pequeno.

A charneira jurídica do Regime~B, para o acoplamento monetário, é o Art.~12 da
Resolução CADE nº~21/2018 \citep{res21cade}, que autoriza considerar o ressarcimento
das vítimas como circunstância atenuante no cálculo da contribuição pecuniária do
TCC (Art.~85 da Lei nº~12.529/2011). O texto vigente do Art.~12 é o seguinte:

\begin{quote}\small
\textit{``Art. 12. A vantagem auferida ou pretendida e o efetivo prejuízo causado
serão considerados, entre outros fatores, na fixação do valor da contribuição
pecuniária, observado o disposto no art. 85 da Lei nº~12.529, de 2011. \S~1º Será
considerada como circunstância atenuante, para fins do disposto no caput deste
artigo, o efetivo ressarcimento extrajudicial ou judicial das vítimas pelo
representado, na forma do art. 45, V e VI da Lei nº~12.529, de 2011.''}
\end{quote}

A recompensa paga pela firma aos denunciantes \emph{pode ser} re-caracterizada como
ressarcimento extrajudicial sob este dispositivo --- e essa re-caracterização é
juridicamente controvertida, sujeita a validação posterior pelo Judiciário
---, gerando o desconto $D$ sem mudança legal. O Regime~C estende
a Lei nº~13.608/2018 (com a redação da Lei nº~13.964/2019) \citep{lei13608} ao
enforcement antitruste, com percentual explícito de recompensa --- maior robustez
jurídica ao custo de via legislativa. O Regime~A é a situação atual, sem canal de
incentivo individual. Crucialmente: o \emph{canal} em si dispensa o Art.~12, com
base autônoma no Art.~4º da Lei nº~12.529/2011 \citep{lei12529} c/c a Lei
nº~9.784/1999 \citep{lei9784}; o Art.~12 é necessário apenas para o acoplamento
monetário-recompensa.

A dinâmica de coordenação da fase P2.5 é ilustrada na Figura~\ref{fig:fase}: a
probabilidade de cascata cresce com a severidade percebida $\sigma$ e cai com o
limiar relativo $k/n$, definindo uma fronteira entre silêncio e revelação.

\begin{figure}[t]
  \centering
  \includegraphics[width=0.78\textwidth]{02_fase_jogo_global}
  \caption{Diagrama de fase da coordenação intrafirma (heurística inspirada em
  \citealp{morris1998unique}). A região verde indica alta probabilidade de cascata
  ($P(\sum a_i \geq k)\to 1$); as estrelas marcam os pontos-alvo dos Regimes B e C.
  \emph{O modelo computacional usa limiares heurísticos, não o equilíbrio fechado
  do jogo global --- ver Limitações.}}
  \label{fig:fase}
\end{figure}

\section{Modelo baseado em agentes e teste de estresse}
\label{sec:limitacoes}

A análise empírica da LCMC usa modelo baseado em agentes (\emph{agent-based
model}, ABM) com três classes de agente --- \texttt{TrabalhadorAgent},
\texttt{EmpresaAgent}, \texttt{AutoridadeAgent} --- implementado em Mesa~3.x.
O \texttt{TrabalhadorAgent} observa conduta com probabilidade dependente do
papel (gradiente Near-Miceli em 3 níveis: primário 1,0; adjacente 0,5;
distal 0,1) e decide depositar segundo seis arquétipos comportamentais
(ético, imitativo, racional, fairminded, aleatório, oportunista) inspirados
em \citet{hokamp2010income}. O \texttt{EmpresaAgent} carrega cultura de
conformidade, conduta potencial específica, e decide cooperação com canal
quando a aritmética da IC-F\* fecha. O \texttt{AutoridadeAgent} mantém o
escrow de depósitos condicionais e processa casos sob restrição de
capacidade $\kappa$ no espírito de \citet{harrington2015when}.

A integração das três populações é dinâmica adaptativa em horizonte
discreto, com seis fases por tique: (P0) atualização de detecção percebida
e ganhos $g_i$ preventivos; (P1) sinalização individual; (P2) depósito
condicional ao escrow; (P2.5) verificação do gatilho de massa crítica
intra-firma e abertura simultânea quando $q_{\min}\cdot n$ é atingido; (P3)
decisão da firma sob IC-F\*; (P4) processamento pela autoridade. A
ortogonalidade entre o canal (P2/P2.5) e os instrumentos de internalização
(P3) é o que sustenta a tese substantiva de que LCMC é mecanismo de
coordenação, com instrumentos como acoplamentos opcionais. A
operacionalização do gatilho de massa crítica é \emph{inspirada} no
limiar de \emph{switching} do jogo global de \citet{morris1998unique} mas
não derivada formalmente dele --- o ABM usa heurística operacional sobre
a fração $q_{\min}\cdot n$ de cooperadores, não o equilíbrio fechado em
$\tau\to 0$.

A calibração formal segue duas dimensões. O gradiente de descontos
intra-firma e inter-firma é calibrado verbatim contra \citet{saito2021tcc}
(349 TCCs do CADE 2012--2019, médias por posição SG/CADE conferidas a partir
da fonte primária). A capacidade institucional é calibrada contra os
Relatórios Integrados de Gestão TCU/CADE 2022--2024 (180 servidores na
área-fim, 73 investigações instauradas em 2024, 109 leniências acumuladas em
20 anos). Os demais parâmetros do modelo seguem ordens de grandeza
plausíveis com pendência de calibração formal.

\subsection{A predição \texorpdfstring{$N^\star$}{N*} e seu teste de falsificação}
\label{sec:nstar}

A calibração formal ajusta o modelo a um alvo operacional de fluxo --- TCCs
por ano --- normalizado pelo número de firmas simuladas. Como o alvo empírico
(a ordem de 47 TCCs/ano observada por \citet{saito2021tcc} sobre o universo
inteiro sob jurisdição do CADE) refere-se a um universo muito maior que as
20 firmas simuladas, a razão de normalização implica um \emph{universo latente}
$N^\star$ de firmas para o qual a taxa simulada seria consistente com a taxa
observada. O ponto ótimo entrega $N^\star \approx 1\,679$.

Trata-se de uma predição derivada, não de um dado ajustado: $N^\star$ não
entra na função objetivo, mas dela decorre. Seu valor é, portanto,
\emph{falsificável} --- e o teste de falsificação é preciso. Cruze-se
$N^\star$ com o número real de firmas sob jurisdição ativa do CADE,
operacionalizado como firmas acima do limiar de faturamento que a Lei
nº~12.529/2011 (Art.~88) fixa para notificação obrigatória de atos de
concentração (R\$~75 milhões para uma das partes), filtradas pelas divisões
CNAE de tecnologia da informação e serviços de informação (62 e 63). Se o
universo real ficar dentro da mesma ordem de grandeza de $N^\star$ --- entre
algumas centenas e alguns milhares ---, a calibração da capacidade
institucional é consistente; se ficar uma ordem de grandeza fora (abaixo de
$\sim$200 ou acima de $\sim$15.000), o parâmetro \emph{taxa\_capacidade}
requer revisão substantiva.

Uma verificação de sanidade parcial já é possível com os dados que
sustentam a calibração de capacidade. A Superintendência-Geral instaurou 73
investigações em 2024 (RIG/TCU-CADE); sob $N^\star \approx 1\,679$, isso
corresponde a uma taxa anual de instauração de $73/1\,679 \approx 4{,}3\%$ do
universo latente --- uma cobertura investigativa da ordem de grandeza
esperada para uma autoridade antitruste de capacidade limitada, nem
negligenciável ($<0{,}5\%$) nem implausivelmente alta ($>30\%$). A predição
sobrevive, portanto, ao único teste que os dados hoje disponíveis permitem;
o teste definitivo contra o cadastro de jurisdicionados fica registrado como
pendência empírica explícita, não realizada aqui para não substituir dado
verificável por estimativa.

\subsection{Proposições}

\begin{proposition}[Viabilidade IC do caminho ``empresa paga'']
Sob Regime B com acoplamento monetário ativo e canal de depósito condicional
operante, existem parâmetros no interior do espaço factível em que IC-F* ($D > W$)
e IR-W são satisfeitas estritamente.
\end{proposition}

\begin{proposition}[Unicidade do equilíbrio de coordenação]
No limite $\tau \to 0$ do subjogo de jogo global, há equilíbrio único de
\textit{switching} $s^*$ para cada $(k, W, r)$ na região relevante; sob
heterogeneidade de arquétipos e papéis, a unicidade é conjectura aberta.
\end{proposition}

\begin{proposition}[Dominância de bem-estar do Regime B]
Para um conjunto de medida positiva de $(W, D, \sigma)$, o bem-estar social
esperado é estritamente maior sob Regime B do que sob Regime A.
\end{proposition}

\begin{proposition}[Coordenação via canal de depósito condicional (LCMC)]
Para qualquer regime de internalização (B ou C), a probabilidade de uma firma com
violação real ser objeto de caso instaurado é estritamente maior sob canal de
depósito condicional com abertura simultânea do que sob denúncias individuais
isoladas, fixados $q_{\min}$, $n$ e a distribuição de arquétipos. Verificada por
construção: a abertura simultânea elimina o equilíbrio de silêncio induzido pelo
custo individual de represália suportado em isolamento.
\end{proposition}

\begin{proposition}[Erosão Coleman --- forma fraca]
A instrumentalização da denúncia interna erode o substrato cooperativo
intra-firma: o capital social residual decai monotonicamente em
$\alpha_{\text{erosão}}$. A operacionalização do parâmetro
\texttt{alpha\_erosao} sobre a fração de notificações por ciclo
\citep{coleman1990foundations, ostrom1990governing} permite verificar
a hipótese empiricamente.
\end{proposition}

\paragraph{Refutação numérica da forma forte.} A formulação inicial da
Proposição~5 (``existe $\alpha^*$ tal que, para $\alpha > \alpha^*$, o
Regime~B colapsa em A''; rastreada em R26) foi falsificada por varredura
em grade (10 seeds $\times$ 8 valores de $\alpha \in [0, 0{,}9]$ $\times$
40 ciclos): a mediana do dano em B fica em $\sim 53$ contra piso A em
$\sim 425$ (fator 8$\times$ menor), estável até $\alpha = 0{,}9$
(Figura~\ref{fig:erosao}). A dissuasão endógena compensa a erosão
do substrato no nível agregado. A Proposição~5 é portanto rebaixada para
a forma fraca acima; a forma forte fica registrada como conjectura
refutada na configuração testada.

\begin{figure}[t]
  \centering
  \includegraphics[width=\textwidth]{03_alpha_erosao_limiar}
  \caption{Falsificação numérica da forma forte da Proposição~5
  (10 seeds $\times$ 8 valores de $\alpha_{\text{erosão}}$ $\times$ 40
  ciclos; bandas = IC bootstrap 95\%). \emph{Esquerda:} dano acumulado em
  Regime~B por $\alpha$ contra o piso do Regime~A (tracejado) --- sem
  cruzamento até $\alpha = 0{,}9$. \emph{Direita:} o capital social
  residual decai monotonicamente --- a forma fraca (erosão do substrato)
  se confirma.}
  \label{fig:erosao}
\end{figure}

\begin{figure}[t]
  \centering
  \includegraphics[width=\textwidth]{04_sankey_lcmc}
  \caption{Fluxo agregado da corrida LCMC sob canal de depósito
  condicional explícito (cenário canônico, 10 firmas $\times$ 120
  trabalhadores $\times$ 15 ciclos): sinais $\to$ depósitos em escrow
  $\to$ aberturas simultâneas $\to$ TCCs assinados. Evidência
  construtiva da Proposição~4.}
  \label{fig:sankey}
\end{figure}

\paragraph{Status (v0.2.0).} A Proposição~1 é verificada por teste de regressão
no ponto-alvo ($D > W$). A Proposição~2 ganhou figura computacional dedicada
(multiplicidade × unicidade Morris-Shin); a unicidade sob heterogeneidade
segue conjectura aberta. A Proposição~3 segue conjectura, sustentada por
varredura de Sobol direcional e por intervalo de confiança bootstrap multi-seed
que não cruza zero. A Proposição~4 é verificada por construção do canal
(abertura simultânea; Figura~\ref{fig:sankey}) e exercitada por cenário
canônico \texttt{apenas\_canal\_sem\_instrumento}. A Proposição~5 tem a forma
forte \emph{refutada} pela varredura dedicada (Figura~\ref{fig:erosao}) e a
forma fraca verificada; o que pode reverter o veredicto --- erosão propagada
por rede inter-firma e calibração de $\alpha$ contra dados reais --- está
rastreado em R26. Os itens de pesquisa que sustentariam plenamente cada
proposição estão rastreados no backlog do repositório.

\section{Enquadramento cibernético}
\label{sec:cibernetica}

\subsection{A Lei da Variedade Requisitada}

A Lei da Variedade Requisitada, enunciada por \citet{ashby1956introduction} em
\emph{An Introduction to Cybernetics}, estabelece que apenas variedade pode
absorver variedade: todo regulador eficaz de um sistema deve conter, em sua
organização interna, repertório de estados distinguíveis pelo menos compatível
com o repertório do sistema regulado. Variedade, no vocabulário ashbyano, é a
cardinalidade do conjunto de estados que um sistema pode ocupar; regular
significa selecionar, na intersecção entre os estados acessíveis e os estados
admissíveis, um subconjunto compatível com o critério normativo do regulador.
Quando a variedade interna do regulador é estritamente menor que a do regulado,
parte do espaço de estados escapa por necessidade estrutural à supervisão ---
não por falha de diligência, mas por restrição informacional intrínseca ao
acoplamento entre os dois sistemas.

\subsection{O teorema de Conant-Ashby}

\citet{conant1970good} reforçam a Lei da Variedade Requisitada com proposição
mais exigente: \emph{todo regulador eficaz de um sistema deve ser um modelo
desse sistema}. A eficácia regulatória não decorre apenas de capacidade
computacional bruta, mas da existência, interna ao regulador, de uma
representação isomorfa --- no nível relevante --- ao sistema regulado. Sob esta
lente, qualquer assimetria entre a estrutura cognitiva do regulador e a
estrutura operacional do regulado define um teto à supervisão factível,
independentemente do volume de recursos materiais alocados.

\subsection{Aplicação ao antitruste digital}

A firma de plataforma contemporânea opera centenas de produtos interdependentes,
milhares de subsistemas algorítmicos de ranqueamento, recomendação e leilão
publicitário, dezenas de modalidades distintas de monetização e aquisição
(\emph{corp dev}, \emph{killer acquisitions}, integração vertical via API
proprietária). Cada subsistema gera estados operacionais distinguíveis do ponto
de vista da licitude antitruste mas internamente opacos para qualquer
observador externo. A área-fim da Superintendência Geral do CADE contava, em
2024, com aproximadamente 180 servidores (RIG/TCU 2024). A defasagem entre a
variedade interna da firma regulada e a variedade investigativa da autoridade
é, portanto, de várias ordens de magnitude. Esta defasagem fornece explicação
estrutural --- e não psicológica nem político-institucional --- para a seleção
subdeterminada das condutas instauradas: o universo de condutas potencialmente
investigáveis excede de tal forma a capacidade do regulador que toda escolha
agregada aproxima-se de uma amostragem com forte componente discricionário.

\subsection{A LCMC como amplificador de variedade}

A LCMC, lida ciberneticamente, não substitui o regulador externo nem amplia seu
efetivo --- amplifica a variedade efetiva do sistema regulatório por recrutamento
de sensores internos à firma. Os trabalhadores que operam o produto, calibram o
ranking, executam o \emph{corp dev} ou desenham a política de monetização
dispõem, por construção da divisão operacional do trabalho, de representação
interna do subsistema relevante equivalente ou superior à do regulador externo.
O canal de depósito condicional realiza a operação que \citet{beer1972brain},
no \emph{Brain of the Firm}, designa de \emph{System~3} no Modelo de Sistemas
Viáveis: estrutura de controle operacional que importa variedade do
\emph{System~1} (operação cotidiana) sem destruir a autonomia subsistêmica.
A importação ocorre por isomorfismo parcial e seletivo, sem captura integral do
estado operacional --- preserva a integridade funcional da unidade observada ao
mesmo tempo em que injeta seu repertório de estados no aparato regulatório.

\subsection{Escrow e preservação da autonomia individual}

O escrow é o mecanismo concreto que torna a amplificação compatível com a
autonomia do trabalhador. O depósito isolado não dispara procedimento --- preserva
o trabalhador da exposição prematura e da represália individualizada que, em
contexto de leniência clássica, sinaliza a defecção a co-conspiradores
identificáveis. A abertura simultânea, condicionada ao gatilho $q_{\min}\cdot n$,
é o momento em que a variedade interna acumulada no escrow é convertida em
variedade disponível ao regulador. A condição $q_{\min}\cdot n$ não é cota
arbitrária: corresponde ao limiar abaixo do qual a importação de variedade seria
insuficiente para distinguir conduta sistêmica de incidente isolado, e acima do
qual o sinal agregado torna-se qualitativamente distinto da soma das observações
individuais.

\subsection{Implicação metodológica}

A leitura cibernética justifica por que a falha da leniência clássica em
mercados digitais não exige hipóteses sobre captura, politização ou má-fé da
autoridade. A defasagem de variedade é estrutural; explica a sub-iniciação sem
recurso a explicações motivacionais ou político-institucionais. A LCMC
posiciona-se, em consequência, como mecanismo de \emph{amplificação} da
variedade investigativa --- não de substituição da autoridade. Esta distinção é
analiticamente relevante: amplificar variedade preserva a alocação
constitucional de competência entre regulador, regulado e Judiciário, ao passo
que substituí-la implicaria reformulação institucional cuja viabilidade
política é distinta da viabilidade técnica do mecanismo aqui proposto.

\section{Pré-condição sociológica}
\label{sec:sociologia}

\subsection{Conspirador externo e trabalhador regular}

A leniência clássica opera sobre um sujeito sociológico bem determinado: o
conspirador externo, vinculado à firma por relação contratual cooperativa de
natureza criminosa, cuja cooperação com a autoridade é tecnicamente uma defecção
dentro de um dilema do prisioneiro pré-existente entre concorrentes. O cálculo
é de cooperação assimétrica sob risco mútuo, e o aparato de incentivos exige
apenas que a imunidade seja preferível, no nó de decisão individual, à
manutenção da conspiração. A LCMC desloca-se para um sujeito qualitativamente
distinto: o trabalhador regular da firma, cuja relação contratual é trabalhista
convencional e cuja decisão de depositar não pertence ao registro da defecção
em conspiração, mas ao registro do abandono parcial da lealdade institucional
em favor de norma externa. O conteúdo psicológico, o substrato normativo e o
horizonte de represália são distintos em natureza, não apenas em grau --- e o
desenho institucional precisa acomodá-los como tais. O depósito não é confissão
de cumplicidade; é exercício de dissenso normativo no interior de relação
contratual cuja conformidade legal é, em princípio, presumida.

\subsection{O fenômeno empírico 2018--2024}

A janela 2018--2024 oferece evidência observacional de que o sujeito requerido
pela LCMC existe empiricamente em empresas de plataforma. O Google Walkout, em
novembro de 2018, mobilizou aproximadamente vinte mil trabalhadores em mais de
cinquenta escritórios em protesto contra a prática de acordos extrajudiciais
sigilosos em casos de assédio --- exercício de ação coletiva sem recompensa
monetária associada e sob proteção jurídica limitada. A constituição da
\emph{Alphabet Workers Union} em janeiro de 2021 é o primeiro caso documentado
de sindicalização em escala relevante numa firma do núcleo da indústria de
plataforma. As denúncias internas de Frances Haugen sobre o Facebook/Meta em
2021, conduzidas sob a proteção restrita do \emph{SEC Whistleblower Office} e
com vazamento subsequente de documentos internos, completam o quadro. Cada
episódio exibe padrão recorrente: denúncia interna coletiva motivada
principalmente por consideração ética, sob proteção legal limitada e sem
contrapartida pecuniária estruturada. O padrão observado é tanto evidência da
existência do sujeito quanto diagnóstico da insuficiência do instrumental
institucional disponível para canalizar sua disposição.

\subsection{Olson e a faixa numérica relevante}

\citet{olson1965logic}, em \emph{The Logic of Collective Action}, previu o
fenômeno da ação coletiva por motivação ética como ocorrência marginal,
restrita a grupos pequenos em que o interesse comum seja suficientemente
perceptível para superar o cálculo do carona. A transposição ao contexto das
plataformas digitais contemporâneas sugere que o número absoluto de
trabalhadores eticamente motivados por firma é pequeno em termos relativos ---
da ordem de cinco a vinte por firma de porte relevante --- mas estritamente
não-nulo. Esta é precisamente a faixa numérica que o gatilho $q_{\min}\cdot n$
da LCMC pretende mobilizar. O canal de depósito condicional não pressupõe
mobilização majoritária; pressupõe a existência de uma minoria eticamente
motivada cuja coordenação simultânea é impedida não pela escassez do
contingente, mas pela ausência de instrumento que sincronize a revelação. O
ponto olsoniano é menos uma profecia pessimista sobre ação coletiva do que uma
cartografia das condições estruturais sob as quais ela ocorre --- e o escrow
é desenho institucional que torna parte dessas condições endógena.

\subsection{Substrato normativo trabalhista}

A LCMC se beneficia de substrato normativo pré-existente, sem dele depender
exclusivamente. Nos Estados Unidos, o \emph{National Labor Relations Board}
reconhece, sob a \emph{Section~7 protected activity} do \emph{National Labor
Relations Act}, a proteção da ação coletiva concertada para mútuo apoio em
condições de trabalho --- cobertura que se estende, doutrinariamente, a formas
de denúncia interna concertada. No Brasil, o Art.~543 da Consolidação das Leis
do Trabalho protege o dirigente sindical contra retaliação direta, e o Art.~462,
\S2º, veda o desconto salarial por ato de cidadania, o que inclui o exercício de
denúncia a órgão público. Estes dispositivos não bastam para induzir a denúncia
em isolamento --- a represália informal, a marginalização interna e a sanção
reputacional ficam fora de sua cobertura ---, mas constituem substrato sobre o
qual o escrow opera como amplificador, sincronizando revelações que, isoladas,
seriam absorvidas pela retaliação individualizada. A LCMC não requer reforma
trabalhista prévia: incorpora a proteção horizontal existente e a complementa
com sincronia temporal.

\subsection{Implicação para Hirschman}

O quadro de \citet{hirschman1970exit} articula três respostas ao declínio
institucional: saída (\emph{exit}), voz (\emph{voice}) e lealdade
(\emph{loyalty}). Em firmas de plataforma com custos elevados de transição
profissional e forte captura reputacional pelo empregador, o trade-off entre
voz e saída tende a deslocar o trabalhador para o silêncio ou para a saída
prematura sem voz. O depósito condicional configura uma forma de voz coletiva
cujo custo individual é diferido: o trabalhador deposita sem se expor, e a
exposição só ocorre quando coordenada com co-depositantes da mesma firma. Em
consequência, a LCMC eleva a viabilidade da opção de voz sem exigir o custo
prévio de saída nem o sacrifício individual da lealdade institucional. Esta
reconfiguração do trade-off voz-saída é, sociologicamente, o que distingue o
depósito condicional da denúncia tradicional: o canal não pede que o trabalhador
escolha entre fidelidade e protesto, mas viabiliza um protesto cuja própria
condição de eficácia é a sincronia com pares motivados pelo mesmo registro
normativo.

\section{Tradução política para o Brasil}
\label{sec:traducao_politica}

A implementação da LCMC no Brasil enfrenta duas portas regulatórias distintas.
A porta infralegal --- Regime B --- consiste em Resolução complementar à
Resolução CADE n.º 21/2018 que institucionalize o procedimento administrativo
de depósito condicional como recepção qualificada de denúncia ao amparo do
Art.~4º (II e III) da Lei 12.529/2011, combinado com a Lei 9.784/99 sobre
processo administrativo federal. A porta legal --- Regime~C --- consiste em
extensão da Lei n.º 13.608/2018 (com a redação da Lei n.º 13.964/2019) ao
enforcement antitruste, criando percentual estatutário explícito de
recompensa ao denunciante.

A análise da viabilidade política 2024--2027 sugere assimetria notável entre
as duas portas. O Regime~B depende exclusivamente de decisão administrativa
do CADE; sua viabilidade é função (i) da prioridade interna do tema na
agenda do Conselho, (ii) da capacidade institucional do DEE e SG para
operacionalizar o canal, e (iii) da aceitação judicial da
re-caracterização do pagamento como ressarcimento sob o Art.~12 da Resolução
21/2018 (risco de anulação judicial, \S\ref{sec:limitacoes}). Nenhuma destas
três condições requer alteração legislativa; todas estão sob arbítrio
administrativo do CADE ou de revisão judicial subsequente.

O Regime~C, em contraste, requer aprovação congressual em um horizonte
político fragmentado: o Congresso brasileiro 2024--2027 tem agenda
concentrada em reforma tributária e PL 2768/2022 (telecomunicações)
está parado desde 2023; o cenário base é de baixa probabilidade de
tramitação de projeto novo de antitruste estatutário. A análise detalhada
em \texttt{viabilidade\_regime\_c.md} do repositório associado conclui que
o Regime~C é \emph{provavelmente infactível} no horizonte 2024--2027 salvo
crise reputacional grande envolvendo plataformas digitais e CADE.

A coordenação inter-institucional CADE--CGU--MPF é segunda dimensão da
tradução política. O canal LCMC do CADE coexistiria com o canal de
denúncia anticorrupção da CGU sob a Lei 12.846/2013 (LAC) e com a
persecução criminal sob a Lei 8.137 conduzida pelo MPF; cada canal
opera competência distinta e proteção distinta ao depositante. A
proposta deste projeto é \emph{não} colocá-los em concorrência, mas
preservar especialização por matéria, com mecanismo de cross-referral
(o trabalhador que deposita matéria de competência mista é encaminhado
ao canal de competência primária).

A comparação com a trajetória DOJ-ATR \citep{dojatr2025} é instrutiva. O
Departamento de Justiça americano instituiu seu programa de recompensa
antitruste em julho de 2025 sob a forma de Resolução administrativa em
parceria com o USPS --- exatamente a estrutura institucional do
Regime~B brasileiro proposto: ato infralegal sob autoridade preexistente,
sem dependência de nova lei. O sucesso (primeiro prêmio pago em janeiro
de 2026) sugere que a via infralegal é viável institucionalmente quando
adotada com decisão administrativa firme. A DMA Whistleblower Tool da
Comissão Europeia \citep{dmatool2024} é precedente complementar para a
via canal-sem-recompensa --- particularmente relevante se a
re-caracterização do Art.~12 enfrentar resistência judicial.

\section{Conclusão}

\subsection{Cinco resultados sustentando a proposta}

Este artigo propôs a Leniência Condicionada à Massa Crítica (LCMC) como
mecanismo de \emph{coordenação} --- não de pagamento --- para o enforcement
antitruste de mercados digitais brasileiros. Cinco resultados sustentam a
proposta. Primeiro, a leniência clássica é estruturalmente inadequada para
condutas unilaterais, predominantes em plataformas digitais: a tecnologia de
ruptura do dilema do prisioneiro entre conspiradores externos perde tração onde
o conhecimento relevante reside numa única firma. Segundo, o canal de depósito
condicional --- \emph{information escrow} no sentido de
\citet{ayres2012information} --- resolve por construção o problema da
sub-iniciação coletiva: a abertura simultânea condicionada ao gatilho
$q_{\min}\cdot n$ elimina o equilíbrio de silêncio induzido pelo custo
individual de represália suportado em isolamento, sem requerer recompensa
monetária prévia. Terceiro, sobre o canal articulam-se cinco instrumentos
opcionais de internalização --- recompensa via TCC, êxodo coletivo
hirschmaniano, crédito tributário, leniência criminal individual, e
intervenção sobre cláusulas restritivas de mobilidade ---, cada qual com
reserva constitucional distinta e analisável separadamente, sem alterar a
estrutura do canal. Quarto, a implementação infralegal sob Regime~B é viável
\emph{de lege lata} pela base autônoma do Art.~4º, II e III, da Lei 12.529
combinado com a Lei 9.784/99, sem dependência da re-caracterização
controvertida do Art.~12 da Resolução 21/2018; o paralelo institucional do
\citet{dojatr2025} (julho de 2025) confirma a viabilidade da via administrativa
em jurisdição comparável. Quinto, o modelo baseado em agentes identifica cinco
vetores de quebra mensuráveis (TCC clássico absorvente; anulação judicial; custo legal individual; massa crítica inalcançável; erosão Coleman do
substrato cooperativo), cada um com parâmetro falsificável e teste de regressão
associado --- a forma forte da Proposição~5 foi explicitamente \emph{refutada}
pela varredura multi-seed dedicada (10 seeds $\times$ 8 valores de $\alpha$
$\times$ 40 ciclos; Figura~\ref{fig:erosao}): o dano em Regime~B fica em
mediana cerca de oito vezes abaixo do piso estável do Regime~A até $\alpha =
0{,}9$, e a dissuasão endógena compensa, no nível agregado, a erosão do
substrato local.

\subsection{Limites reconhecidos}

Os limites da contribuição são reconhecidos. A calibração formal
identifica $N^\star \approx 1\,679$ firmas como predição falsificável
(\S\ref{sec:nstar}); a predição sobrevive ao teste de sanidade que os dados
públicos permitem --- taxa de instauração de $\approx 4{,}3\%$ do universo
latente ---, mas o teste definitivo contra o cadastro de jurisdicionados do
CADE permanece pendência empírica não realizada aqui. A
Proposição~2, no que se refere à unicidade do equilíbrio de \emph{switching}
sob heterogeneidade de arquétipos e papéis, permanece conjectura aberta ---
verificada por construção do canal no limite homogêneo $\tau \to 0$, mas
formalmente não fechada no caso heterogêneo. A viabilidade
política do Regime~C no horizonte 2024--2027 é baixa, dada a agenda congressual
concentrada em reforma tributária e a paralisia do PL 2768/2022; a
viabilidade do Regime~B, ainda que tecnicamente preparada, depende de decisão
administrativa do CADE não anunciada nem em consulta pública. A análise
comparada com programas existentes sugere que o desenho ocupa nicho ainda
vazio na taxonomia institucional contemporânea --- canal de depósito
condicional com massa crítica intra-firma sob via infralegal ---, de modo que
a contribuição substantiva é localizar e formalizar este nicho, não demonstrar
superioridade incondicional sobre alternativas em qualquer cenário.

\subsection{Três frentes futuras}

A análise futura deve perseguir três frentes. A frente empírica consiste em
confirmar ou refutar $N^\star$ contra cadastro CADE de firmas sob jurisdição
ativa; e em calibrar $\alpha_{\text{erosão}}$ contra evidência longitudinal de
denúncia em jurisdições com programas estabelecidos --- séries históricas da
SEC, dados a serem produzidos pelo DOJ-ATR a partir de 2026, e métricas da
DMA Whistleblower Tool. A frente teórica consiste em fechar a unicidade do
equilíbrio sob heterogeneidade (Proposição~2 sob a forma heterogênea) e em
investigar formalmente as condições de não-bifurcação no acoplamento
$(\lambda, \beta_H)$ entre laços dissuasório e contratual via análise de
jacobiano local. A frente institucional consiste em redigir, em sede
acadêmica colaborativa com pesquisadores de direito administrativo, minuta de
Resolução complementar à Resolução CADE n.º 21/2018 que articule procedimento
formal de depósito condicional --- contribuição que este artigo declina por
não-vinculação política do autor, mas que fica registrada como pendência
metodologicamente preparada para trabalho coordenado subsequente.

\bibliographystyle{apalike}
\bibliography{refs}

\end{document}
